\chapter{Finite State Automata and Regular Languages}

\section{Introduction}

In this first chapter, we introduce the fundamental concepts of automata theory, which serves as the main tool for representing systems in verification and validation. Systems are typically described by means of a set of states and a set of transitions. 

A key assumption in this course is that we will be dealing with \textbf{Finite State Automata}. A finite state automaton is an automaton where the number of states (or nodes) is strictly finite. 

Alongside finite state automata, systems can also be represented using \textbf{Regular Expressions}. Throughout the course, we will explore both formalisms and establish the correspondence between them.

\section{Alphabets, Symbols, and Words}

The basic building block for our definitions is a \textbf{set of symbols} (often denoted by letters like A, B, C or binary digits 0, 1). We refer to a finite set of symbols as an \textbf{alphabet}, typically denoted by $\mathcal{A}$.

Using symbols, we can build \textbf{words}. A word is simply a sequence of symbols from the alphabet. For example, if the alphabet is $\{A, B, C\}$, a word could be $ABCAB$. In this preliminary part of the course, we focus on \textbf{finite words}, meaning the length of the sequence is finite. 

\begin{note}
Later in the course, we will transition to infinite words to model systems that conceptually run forever (like operating systems) and do not have a defined termination point.
\end{note}

We commonly use variables like $u, v, w$ to denote words.

\subsection{Length of a Word}

The length of a word $w$, denoted by $|w|$, is the number of symbols it contains. For example, if $w = ABCDC$, then $|w| = 5$.

A special word in formal languages is the \textbf{empty word}, denoted by $\varepsilon$. The empty word contains no symbols, meaning its length is zero: $|\varepsilon| = 0$.

\subsection{Operations on Words: Concatenation}

The primary operation on words is \textbf{concatenation}. Given two words $u$ and $v$, their concatenation, denoted as $u \cdot v$ (or simply $uv$), is the sequence of symbols of $u$ immediately followed by the sequence of symbols of $v$. 

Properties of concatenation:
\begin{itemize}
\item \textbf{Length}: $|uv| = |u| + |v|$
\item \textbf{Identity element}: Concatenating any word $u$ with the empty word yields $u$. That is, $u\varepsilon = \varepsilon u = u$.
\item \textbf{Non-commutative}: In general, $uv \neq vu$ (except in special cases, like $u = AA$ and $v = AAA$).
\end{itemize}

\section{Operations on Sets of Words}

Concatenation can be naturally extended to sets of words. If $A$ and $B$ are sets of words, their concatenation $AB$ is defined as:

$$ AB = \{ uv \mid u \in A \text{ and } v \in B \} $$

\textbf{Important edge cases}:
\begin{itemize}
\item If $A$ is concatenated with $\{\varepsilon\}$ (a set containing only the empty word), the result is $A$.
\item If $A$ is concatenated with the empty set $\emptyset$, the result is $\emptyset$ (since you cannot pick a word from an empty set).
\end{itemize}

\subsection{The Kleene Star (Closure)}

We can define the $n$-th power of a set $A$, denoted $A^n$, as the set obtained by concatenating $A$ with itself $n$ times.
\begin{itemize}
\item $A^0 = \{\varepsilon\}$ (by convention)
\item $A^1 = A$
\item $A^{n+1} = A \cdot A^n$
\end{itemize}

The \textbf{Kleene Closure} (or Kleene Star) of a set $A$, denoted $A^*$, is defined as the infinite union of all its powers:

$$ A^* = \bigcup_{n \ge 0} A^n $$

Essentially, $A^*$ is the set of all possible finite words that can be constructed by concatenating an arbitrary (but finite) number of words from $A$. 

\begin{note}
A related operation is $A^+$, which is defined as $A^* \setminus \{\varepsilon\}$, assuming $\varepsilon \notin A$.
\end{note}

\section{Formal Languages}

A \textbf{language} $L$ over an alphabet $\mathcal{A}$ is defined formally as any subset of $\mathcal{A}^*$.
$$ L \subseteq \mathcal{A}^* $$

Typically, the languages we study will contain an infinite number of words, even though each individual word is of finite length.

\section{Deterministic Finite State Automata (DFA)}

We are now ready to define a \textbf{Deterministic Finite State Automaton (DFA)}. 
A DFA is mathematically represented as a 5-tuple: $\mathcal{A} = (Q, \mathcal{A}, \delta, q_0, F)$
where:
\begin{enumerate}
\item $Q$: A finite set of states.
\item $\mathcal{A}$: A finite set of symbols (the input alphabet).
\item $\delta$: The transition function, $\delta: Q \times \mathcal{A} \rightarrow Q$. It takes the current state and an input symbol and returns the next state.
\item $q_0$: The initial state ($q_0 \in Q$).
\item $F$: A set of final (or accepting) states ($F \subseteq Q$).
\end{enumerate}

A DFA is \textbf{deterministic} because for any given state and any given input symbol, there is \emph{at most one} transition defined. You can reach at most one specific state.

\subsection{Incomplete Automata and Sink States}

Sometimes, an automaton might not have a transition defined for every symbol in every state. Such an automaton is called \textbf{incomplete}. 

Any incomplete automaton can be made complete without changing the language it recognizes by introducing a \textbf{sink state} ($q_S$). 
\begin{itemize}
\item For every missing transition from a state $q$ reading symbol $c$, we add a transition $\delta(q, c) = q_S$.
\item Once the automaton enters the sink state, it loops back to the sink state for every symbol in the alphabet ($\delta(q_S, a) = q_S, \forall a \in \mathcal{A}$).
\item The sink state is \emph{not} a final state.
\end{itemize}

\subsection{Extended Transition Function and Acceptance}

To describe how a DFA processes an entire word (a sequence of symbols), we extend the transition function $\delta$ to a function $\hat{\delta}: Q \times \mathcal{A}^* \rightarrow Q$.
\begin{itemize}
\item \textbf{Base case}: $\hat{\delta}(q, \varepsilon) = q$
\item \textbf{Inductive step}: $\hat{\delta}(q, ua) = \delta(\hat{\delta}(q, u), a)$ for a word $u$ and symbol $a$.
\end{itemize}

\textbf{Acceptance Condition}: 
A word $w$ is accepted by a DFA if and only if processing the word from the initial state leads to a final state:
$$ \hat{\delta}(q_0, w) \in F $$

If the automaton stops because of an undefined transition, or if it finishes processing the word but ends up in a non-final state, the word is \textbf{rejected}.

The \textbf{language recognized (or accepted) by a DFA $\mathcal{A}$}, denoted $L(\mathcal{A})$, is the set of all words accepted by the automaton. Any language that can be recognized by a DFA is called a \textbf{regular language}.

\section{Non-Deterministic Finite State Automata (NFA)}

In contrast to a DFA, a \textbf{Non-Deterministic Finite State Automaton (NFA)} allows for multiple possible transitions from a given state for a single input symbol. 

An NFA is defined as a 5-tuple $(Q, \mathcal{A}, \Delta, q_0, F)$, where the first four components and the last are the same as a DFA, but the transition function $\Delta$ is defined differently:
\begin{itemize}
\item $\Delta: Q \times \mathcal{A} \rightarrow 2^Q$ (where $2^Q$ is the power set of $Q$).
\end{itemize}
Alternatively, $\Delta$ can be seen as a relation $\Delta \subseteq Q \times \mathcal{A} \times Q$.

This means that reading a symbol from a state can lead to a \emph{set} of possible states. Non-determinism is useful to represent systems with choices or uncertainty. 

When generalizing the transition function $\hat{\Delta}$ for an NFA to process words, we must track the entire \emph{set of states} the automaton could possibly be in at any given step.
